% !TEX program = lualatex
% Reglamento del Consejo de Estudiantes de la Facultad de Ciencias (Universidad de Zaragoza).
% Compilar con LuaLaTeX (latexmk lo toma de la primera línea; en Overleaf, Menu > Compiler).
% La tipografía Myriad Pro se carga desde la carpeta fonts/.
\documentclass[11pt,a4paper]{article}
\input{estilo-reglamento}
% \borradorfalse   % descomentar para la versión que se firma

\begin{document}

\thispagestyle{portada}
\vspace*{50mm}
\begin{center}
{\fontsize{26}{32}\selectfont\bfseries REGLAMENTO DEL CONSEJO DE\\[2pt] ESTUDIANTES\\[2pt] FACULTAD DE CIENCIAS\par}
\vspace{6pt}
{\fontsize{11}{14}\selectfont DE LA REPRESENTACIÓN DE LOS ESTUDIANTES Y SU ORGANIZACIÓN\par}
\end{center}
\clearpage

\tableofcontents
\clearpage


\capitulo{CAPÍTULO PRIMERO}{De los representantes de estudiantes de la Facultad de Ciencias}

\articulo{1}{Condición de representante}
\begin{enumerate}
\item Ostentan la condición de representantes electos de estudiantes de la Facultad de Ciencias aquellos estudiantes matriculados en el centro que ostenten cualquiera de los siguientes puestos: delegaciones o subdelegaciones de grupos de docencia, miembros del Consejo de Facultad, Claustro Universitario y Consejos de los Departamentos con sede en la Facultad de Ciencias. Ostentan la condición de representantes designados aquellos estudiantes que sean escogidos para representar a los estudiantes en otros órganos colegiados y de representación de la Universidad de Zaragoza o en nombre de la Universidad de Zaragoza, de forma indirecta a través del Pleno del Consejo de Estudiantes.
\item Se perderá la condición de representante electo de estudiantes de la Facultad por:
  \begin{enumerate}
    \item Incumplimiento de las condiciones expuestas en el apartado 1 del presente artículo. A estos efectos, la condición de estudiante matriculado se mantiene durante el primer periodo de matriculación del curso siguiente, salvo para quienes hayan concluido los estudios de la titulación en la que estaban matriculados.
    \item Dimisión, que será efectiva de manera inmediata tras haberse hecho constar una voluntad explícita en tal sentido, mediante escrito formulado ante el Consejo de Estudiantes y obrante asimismo en la Secretaría de la Facultad.
  \end{enumerate}
\item Se perderá la condición de representante designado de estudiantes de la Facultad por:
  \begin{enumerate}
    \item Transcurso del periodo de su nombramiento, sin perjuicio de la posibilidad de reelección. Salvo que la normativa del órgano correspondiente establezca otro periodo, el nombramiento durará hasta la sesión constitutiva del Pleno siguiente.
    \item Pérdida de la condición de estudiante de la Facultad, en los términos de la letra a) del apartado 2 del presente artículo.
    \item Dimisión, que será efectiva de manera inmediata tras haberse hecho constar una voluntad explícita en tal sentido, mediante escrito formulado ante la secretaría del Consejo de Estudiantes.
    \item Revocación acordada por el Pleno, que se tramitará conforme al procedimiento de la moción de censura previsto en el artículo 7.7, o retirada del cargo de representación de conformidad con la normativa del órgano en el que se ejerza.
  \end{enumerate}
\item Las delegaciones y subdelegaciones de grupo continuarán en funciones hasta la proclamación de quienes les sucedan.
\end{enumerate}

\articulo{2}{De los derechos y deberes de los representantes de estudiantes}
\begin{enumerate}
\item Son derechos de los representantes de estudiantes de la Facultad de Ciencias:
  \begin{enumerate}
    \item Tener la consideración de representativos sus actos y manifestaciones, salvo que conste haber sido realizados a título meramente personal.
    \item Solicitar y recibir del Consejo de Estudiantes, del profesorado y de los órganos de gobierno de la Facultad, la información y el apoyo que precisen para el desempeño de sus funciones.
    \item Ser atendidos por el profesorado y los miembros del personal técnico, de gestión y de administración y servicios para que puedan hacer compatible, en la medida de lo posible, la labor de representación con sus estudios, así como para un mejor desempeño de sus funciones.
    \item Todos aquellos que les reconozca la legislación vigente y otras normas de la Universidad.
  \end{enumerate}
\item Son deberes de los representantes de estudiantes de la Facultad de Ciencias:
  \begin{enumerate}
    \item Participar activamente en el desempeño de las funciones representativas de los estudiantes de la Facultad de Ciencias y, particularmente, en lo que se refiere a la defensa de sus intereses.
    \item Transmitir fidedignamente a quien corresponda las peticiones de sus representados, así como las opiniones de los mismos si así lo solicitaran estos.
    \item Informar a sus representados de las actividades realizadas en el desempeño de su cargo, de los hechos relevantes para la vida universitaria y de todas aquellas de interés para los mismos.
    \item Coordinarse y relacionarse con el resto de los representantes.
    \item Conocer, cumplir y hacer cumplir el presente Reglamento.
    \item Todos aquellos que se deriven del presente Reglamento, los inherentes al ejercicio de su cargo y los demás previstos en la legislación vigente.
  \end{enumerate}
\end{enumerate}

\capitulo{CAPÍTULO SEGUNDO}{Del Consejo de Estudiantes de la Facultad de Ciencias}

\articulo{3}{Naturaleza, composición y funciones}
\begin{enumerate}
\item El Consejo de Estudiantes es el órgano colegiado de deliberación, consulta y representación de los intereses de los estudiantes de la Facultad de Ciencias.
\item Son miembros del Consejo de Estudiantes todos los representantes electos de la Facultad de Ciencias. El Consejo estará presidido por quien ostente la presidencia o, en su caso, la vicepresidencia.
\item Son funciones del Consejo de Estudiantes:
  \begin{enumerate}
    \item Representar a todos los estudiantes de la Facultad.
    \item Elegir, de entre sus miembros, a la presidencia y al resto de la Comisión Permanente.
    \item Recabar y facilitar información a los estudiantes sobre los aspectos académicos y universitarios de interés para los mismos.
    \item Recabar los medios necesarios para el ejercicio de sus funciones.
    \item Impulsar aquellos proyectos que se consideren de interés para los estudiantes de la Facultad, así como promover iniciativas a presentar ante los órganos de gobierno de la Facultad, de la Universidad y de otras instituciones de interés.
    \item Contribuir al desarrollo y calidad de la Facultad, colaborando en la mejora de sus servicios e impulsando programas, actos y otras actividades que redunden en beneficio de esta y de sus estudiantes.
    \item Colaborar con sus medios y recursos con el asociacionismo estudiantil.
    \item Pronunciarse sobre las cuestiones que sean sometidas a su consideración por cualquiera de sus miembros y, en especial, por la Comisión Permanente.
    \item Velar por la aplicación del presente Reglamento en aquellas cuestiones de su competencia.
    \item Aprobar, en desarrollo de este Reglamento, las normas de funcionamiento del Consejo de Estudiantes, que deberán ser ratificadas por el Consejo de Facultad.
    \item Elegir entre sus miembros las respectivas delegaciones de titulación.
    \item Elegir entre sus miembros a sus representantes en el Consejo de Estudiantes de la Universidad de Zaragoza (CEUZ).
    \item Elegir entre sus miembros a sus representantes en asociaciones estudiantiles sectoriales y otros colectivos suprauniversitarios.
    \item Todas aquellas que deriven de la legislación vigente.
  \end{enumerate}
\item El Consejo de Estudiantes actuará constituido en pleno y podrá crear tantas comisiones como estime necesarias para su buen funcionamiento.
\end{enumerate}

\articulo{4}{Pleno del Consejo de Estudiantes}
\begin{enumerate}
\item El Pleno del Consejo de Estudiantes se constituirá en un plazo máximo de diez días naturales siguientes al de la proclamación definitiva de los resultados de las elecciones a delegación y subdelegación de los grupos de docencia de la Facultad. Si dicho plazo coincidiera, en todo o en parte, con el periodo en el que no pueden celebrarse plenos ordinarios conforme al apartado 3, la sesión se celebrará dentro de los diez días naturales siguientes a su finalización.
\item La sesión constitutiva tendrá carácter ordinario y será convocada por la presidencia en funciones del Pleno saliente o, en su defecto y sucesivamente, por la vicepresidencia, la secretaría o el Decanato.
\item El Pleno se reunirá, en sesión ordinaria, al menos una vez al cuatrimestre durante el periodo lectivo. Dichas sesiones serán convocadas por su presidencia con una antelación mínima de cinco días lectivos. No podrán celebrarse plenos ordinarios durante los periodos de exámenes fijados en el calendario académico ni en los quince días naturales anteriores a su inicio.
\item Igualmente, el Pleno podrá reunirse, en sesión extraordinaria, convocada por su presidencia, bien a iniciativa propia, bien a solicitud de una quinta parte de sus miembros o bien de la Comisión Permanente. Dichas sesiones serán convocadas con una antelación mínima de dos días naturales. Solicitada la sesión extraordinaria, la presidencia deberá convocarla en un plazo máximo de diez días naturales; en su defecto, la convocará la vicepresidencia o, en defecto de esta, la secretaría.
\item Los miembros del Pleno serán convocados de forma individual mediante correo electrónico a su dirección personal institucional. La convocatoria deberá incluir siempre el orden del día.
\item Los puntos del orden del día de un pleno ordinario deberán obligatoriamente cubrir las siguientes cuestiones:
  \begin{enumerate}
    \item Intervención, si procede, del equipo decanal.
    \item Aprobación, si procede, de las actas de plenos anteriores.
    \item Rendición de cuentas por parte de los miembros de la Comisión Permanente.
    \item Ruegos y preguntas.
  \end{enumerate}
\item No podrá ser objeto de deliberación o acuerdo ningún asunto que no figure en el orden del día, salvo que asistan todos los miembros del Pleno y sea declarada la urgencia del asunto por mayoría simple.
\item Todos los miembros electos del Pleno tienen derecho a voz y voto en sus sesiones. Cada miembro dispondrá de un único voto propio, con independencia del número de puestos por los que ostente la condición de representante electo.
\item Serán invitados al Pleno con voz, pero sin voto, todos aquellos representantes designados que no fueran ya miembros electos por algún otro cargo.
\item La presidencia podrá invitar a las sesiones del Pleno a cualquier persona que considere de interés, ya sea a iniciativa propia o de una quinta parte de los miembros del Pleno. Las personas invitadas tendrán voz, pero no voto.
\item Para la válida constitución del Pleno será necesaria la asistencia, en primera convocatoria, de al menos un tercio de sus miembros y, en segunda convocatoria, quince minutos después, de al menos un cuarto de sus miembros. Los votos delegados no computarán a efectos de quorum.
\item Cualquier estudiante de la Facultad de Ciencias podrá solicitar a la Comisión Permanente, mediante correo electrónico y hasta dos días naturales antes de la celebración, asistir sin voz ni voto a cualquier pleno ordinario. Todas las solicitudes serán resueltas, como tarde, un día antes de la celebración del pleno, atendiendo a necesidades logísticas y de confidencialidad.
\item Las sesiones del Pleno podrán celebrarse de forma presencial, telemática o mixta. La participación telemática tiene por objeto facilitar intervenciones a distancia y no sustituye la asistencia presencial; quienes participen telemáticamente solo podrán votar en las aprobaciones por asentimiento y en las votaciones que se realicen mediante plataforma en línea conforme al artículo 5. Asimismo, la presidencia podrá someter a aprobación por asentimiento, por vía telemática, cualquier propuesta, salvo la elección de cargos de la Comisión Permanente, la moción de censura y la reforma del Reglamento. La propuesta se enviará a todos los miembros del Pleno por correo electrónico institucional, con un plazo no inferior a dos días naturales, y se entenderá aprobada si al término del plazo ningún miembro se ha opuesto. Si algún miembro se opone, la propuesta se incluirá en el orden del día del siguiente pleno. De los acuerdos así adoptados se dará cuenta en el siguiente pleno ordinario.
\end{enumerate}

\articulo{5}{De las votaciones}
\begin{enumerate}
\item A efectos de este Reglamento:
  \begin{enumerate}
    \item Son votos emitidos los votos a favor, en contra y en blanco. Las abstenciones y los votos no emitidos no computan como votos emitidos.
    \item Hay mayoría simple cuando los votos a favor superan a los votos en contra.
    \item Hay mayoría absoluta cuando los votos a favor superan la mitad de los votos emitidos.
    \item Las mayorías referidas a los miembros del Pleno se computan sobre el número total de sus miembros.
    \item Son miembros presentes quienes asisten a la sesión, presencial o telemáticamente. Quien haya delegado su voto no se considera presente.
  \end{enumerate}
\item Las decisiones se adoptarán, salvo que se especifique lo contrario, por mayoría simple. El quorum previsto en el artículo 4.11 deberá existir en el momento de cada votación.
\item Serán secretas, en todo caso, las votaciones que afecten a personas: elección de cargos, vocalías y representantes designados, mociones de censura y revocaciones. Cualquier otra votación será secreta cuando lo solicite al menos una quinta parte de los miembros presentes. Lo anterior se entiende sin perjuicio de la aprobación por asentimiento prevista en el artículo 4.13.
\item Las votaciones secretas se realizarán, siempre que sea posible, por medios electrónicos que garanticen el anonimato, la unicidad del voto y la verificabilidad del recuento. Cuando no sea posible, se votará mediante papeleta en urna.
\item Salvo lo dispuesto en el artículo 7.1, las votaciones que se realicen mediante plataforma en línea se abrirán al término de la sesión, con el envío del enlace correspondiente, y permanecerán abiertas durante el plazo que anuncie la presidencia antes de levantar la sesión. Solo podrán votar en ellas los miembros presentes en la sesión y, en su caso, quienes ejerzan votos delegados.
\item Cada miembro del Pleno podrá delegar su voto en otro miembro, mediante el procedimiento formal indicado en la convocatoria, justificando su ausencia ante la secretaría. La convocatoria fijará el plazo para comunicar la delegación, que finalizará entre dos días naturales y una hora antes del inicio de la sesión en los plenos ordinarios, y entre veinticuatro horas y una hora antes en los extraordinarios. Ningún miembro podrá ostentar más de dos votos delegados. Los votos delegados computan en el recuento de todas las votaciones, incluidas las secretas y las realizadas en línea, en las que la persona en quien se delega recibirá el enlace de voto de quien delega y votará en su nombre. Los votos delegados no computan a efectos de quorum ni de asistencia.
\end{enumerate}

\capitulo{CAPÍTULO TERCERO}{De la Comisión Permanente del Consejo de Estudiantes}

\articulo{6}{Composición y funciones}
\begin{enumerate}
\item La Comisión Permanente estará compuesta, como mínimo, por:
  \begin{enumerate}
    \item La presidencia del Consejo de Estudiantes.
    \item La vicepresidencia del Consejo de Estudiantes.
    \item La secretaría del Consejo de Estudiantes.
    \item La tesorería del Consejo de Estudiantes.
  \end{enumerate}
\end{enumerate}
Además, la Comisión Permanente contará con las vocalías que se elijan conforme al artículo 7.

\begin{enumerate}[resume]
\item Son funciones de la Comisión Permanente las propias de los artículos 8, 9, 10 y 11 y las demás que le atribuye este Reglamento.
\end{enumerate}

\articulo{7}{Elección y cese}
\begin{enumerate}
\item El Pleno elegirá de entre sus miembros a la Comisión Permanente en la sesión constitutiva prevista en el artículo 4.1. La elección de la presidencia, la vicepresidencia, la secretaría y la tesorería se votará necesariamente durante la propia sesión, de modo que la candidatura en bloque o la presidencia elegidas puedan proponer en ella las vocalías conforme al apartado 4. La votación de las vocalías podrá realizarse mediante plataforma en línea al término de la sesión, conforme al artículo 5.5.
\item Las candidaturas a la presidencia, vicepresidencia, secretaría y tesorería podrán presentarse en bloque o individualmente. Cada candidatura dispondrá de un turno de presentación de al menos un minuto.
\item Si se presenta alguna candidatura en bloque, se votarán en primer lugar las candidaturas en bloque. Si ninguna resultara elegida, o si no se hubiera presentado ninguna, los cargos se elegirán uno a uno, por este orden: presidencia, vicepresidencia, secretaría y tesorería, y podrán concurrir a ellos también quienes integraran candidaturas en bloque. Los cargos que no incluya la candidatura en bloque elegida se elegirán del mismo modo. En cada votación podrá votarse a favor de una candidatura, en blanco o en contra de todas. Si concurre una sola candidatura, resultará elegida si obtiene mayoría simple. Si concurren varias, en primera votación se requerirá mayoría absoluta; si ninguna la obtiene, se celebrará una segunda votación entre las dos más votadas, en la que resultará elegida la que obtenga más votos. Los empates se resolverán repitiendo la votación entre las candidaturas empatadas y, de persistir, por sorteo.
\item Vocalías.
  \begin{enumerate}
    \item Una vez elegidos los cargos del apartado 3, la candidatura en bloque o, en su defecto, la presidencia elegida propondrá las vocalías que estime oportunas, con su denominación y funciones. El Pleno podrá suprimir en cualquier momento una vocalía que haya quedado desierta.
    \item Cualquier miembro del Pleno podrá presentar, durante la sesión, su candidatura a una o varias vocalías, pero solo podrá ocupar una. Quien se presente a varias deberá indicar, al presentar sus candidaturas, su orden de preferencia entre ellas; de no hacerlo, se tomará el orden en que las haya presentado.
    \item Todas las vocalías se votarán de forma simultánea. En cada vocalía, cada miembro del Pleno podrá votar a favor, en contra o abstenerse respecto de cada candidatura. Solo serán elegibles en una vocalía las candidaturas con más votos a favor que en contra en ella.
    \item Las vocalías se asignarán conforme al siguiente procedimiento:
      \begin{ordinales}
        \item Cada candidatura opta a la vocalía que más prefiera de entre aquellas en las que sea elegible y no haya sido descartada.
        \item Cada vocalía retiene provisionalmente, de entre las candidaturas que optan a ella y la que tuviera ya retenida, la que haya obtenido más votos a favor en dicha vocalía, y descarta las demás.
        \item Las candidaturas descartadas optan a su siguiente vocalía preferida y se repite el paso 2.º, hasta que ninguna candidatura sea descartada o las descartadas no puedan optar a ninguna otra vocalía. Las asignaciones provisionales resultantes serán definitivas.
      \end{ordinales}
    \item Los empates en votos a favor dentro de una vocalía se resolverán en favor de la candidatura con menos votos en contra en ella y, de persistir, conforme a un orden único establecido por sorteo antes del recuento.
    \item La vocalía que quede sin titular se declarará desierta y se procederá conforme al apartado 5. Si una vocalía quedara vacante con posterioridad, se ofrecerá a la siguiente candidatura elegible en ella que no ocupe otra vocalía; en su defecto, se procederá conforme al apartado 5.
    \item La secretaría comunicará por correo electrónico a todos los miembros del Pleno los resultados de cada vocalía y la asignación final.
    \item Cuando no se disponga de medios telemáticos que garanticen las condiciones del artículo 5.4, o cuando así lo solicite al menos una quinta parte de los miembros del Pleno, la Comisión Permanente propondrá un procedimiento alternativo de votación, que deberá ser refrendado por el Pleno por mayoría simple antes de celebrarse la votación. Entre otros, podrá proponerse:
      \begin{ordinales}
        \item La votación sucesiva de las vocalías, de una en una y en el orden que fije la presidencia, de modo que quien resulte elegido en una quede excluido de las siguientes.
        \item La votación presencial mediante papeleta en urna, manteniendo el procedimiento de asignación de la letra d).
        \item La designación por asentimiento cuando haya una sola candidatura presentada por vocalía y nadie se oponga.
      \end{ordinales}
  \end{enumerate}
\item Si algún cargo de la Comisión Permanente quedara sin cubrir, su elección se incluirá en el orden del día de cada pleno posterior hasta que se cubra. Mientras tanto, las funciones de la presidencia las ejercerá la vicepresidencia, y las de los demás cargos, quien designe la Comisión Permanente. Si no se hubiera elegido ni presidencia ni vicepresidencia, la Comisión Permanente saliente continuará en funciones; en su defecto, el Decanato convocará el Pleno cuando sea necesario elegir representantes y, al menos, una vez por cuatrimestre para repetir la elección. Nadie podrá ocupar un cargo de la Comisión Permanente sin haber presentado su candidatura.
\item El mandato de la Comisión Permanente dura hasta la sesión constitutiva del Pleno siguiente. Hasta entonces, sus miembros continuarán en funciones aunque hayan perdido la condición de representante electo por haber concluido su mandato como delegación o subdelegación de grupo o haber finalizado sus estudios. Sus miembros podrán ser reelegidos.
\item Al menos un tercio de los miembros del Pleno podrá presentar una moción de censura contra la Comisión Permanente en su conjunto o contra alguno de sus miembros. La moción deberá incluir una candidatura alternativa para cada cargo afectado; si se dirige contra la presidencia o contra la Comisión en su conjunto, al menos para la presidencia. La presidencia o, en su defecto, la vicepresidencia convocará en un plazo máximo de diez días naturales un pleno extraordinario con la moción como único punto; si no lo hicieran, lo convocará la secretaría o, en su defecto, la primera persona firmante de la moción. Para su debate y votación se requerirá la asistencia de al menos la mitad de los miembros del Pleno, y la moción se aprobará por mayoría absoluta de los miembros presentes. Si se aprueba contra la presidencia o contra la Comisión entera, cesará toda la Comisión, la candidatura alternativa ocupará los cargos para los que se haya presentado y los restantes se cubrirán conforme a los apartados 3 a 5; si se aprueba contra otro miembro, la candidatura alternativa ocupará su puesto. Quienes firmen una moción que resulte rechazada no podrán presentar otra durante el mismo cuatrimestre.
\item Los miembros de la Comisión Permanente cesarán, además, por perder la condición de representante electo (artículo 1.2), salvo lo previsto en el apartado 6, o por dimitir del cargo. La vacante se cubrirá conforme al apartado 5 o, en el caso de las vocalías, conforme a la letra f) del apartado 4.
\end{enumerate}

\articulo{8}{De la presidencia del Consejo de Estudiantes}
\begin{enumerate}
\item Son competencias de la presidencia:
  \begin{enumerate}
    \item Representar a todos los estudiantes de la Facultad, dentro y fuera de ella.
    \item Presidir y dirigir las sesiones del Pleno.
    \item Actuar de enlace entre los distintos órganos de representación estudiantil.
    \item Velar por la ejecución de los acuerdos tomados en el Consejo de Estudiantes, así como por el correcto uso de los locales y los medios puestos a disposición del Consejo de Estudiantes.
    \item Garantizar la publicidad de todas las convocatorias electorales que afecten a los estudiantes.
    \item Todas aquellas que le sean atribuidas por el presente Reglamento y por el resto de la legislación aplicable.
  \end{enumerate}
\end{enumerate}

\articulo{9}{De la secretaría del Consejo de Estudiantes}
\begin{enumerate}
\item Son competencias de la secretaría:
  \begin{enumerate}
    \item Efectuar la convocatoria de las sesiones del Pleno por orden de la presidencia, así como las citaciones a los miembros invitados de las mismas.
    \item Levantar acta de las sesiones del Pleno.
    \item Llevar al día el registro de los miembros del Consejo de Estudiantes.
    \item Custodiar las actas, los registros y la documentación de las votaciones del Consejo de Estudiantes.
    \item Todas aquellas que le correspondan conforme a este Reglamento y su normativa de desarrollo.
  \end{enumerate}
\end{enumerate}

\articulo{10}{De la vicepresidencia del Consejo de Estudiantes}
\begin{enumerate}
\item Son competencias de la vicepresidencia:
  \begin{enumerate}
    \item Realizar las labores de la presidencia en caso de que esta no se encuentre disponible.
    \item Auxiliar a la presidencia y a la secretaría cuando estas lo requieran.
    \item Todas aquellas que le correspondan conforme a este Reglamento y su normativa de desarrollo.
  \end{enumerate}
\end{enumerate}

\articulo{11}{De la tesorería del Consejo de Estudiantes}
\begin{enumerate}
\item Son competencias de la tesorería:
  \begin{enumerate}
    \item Llevar al corriente las cuentas y asuntos monetarios del Consejo de Estudiantes.
    \item Elaborar presupuestos para las actividades realizadas por el Consejo de Estudiantes y velar por su correcta ejecución.
    \item Gestionar el presupuesto asignado al Consejo de Estudiantes por el vicerrectorado con competencias en materia de estudiantes y la justificación de los proyectos asociados.
    \item Presentar ante el Pleno, al menos una vez por curso académico, un informe detallado de las cuentas, así como cuando lo solicite cualquier miembro del Pleno o la Comisión Permanente. En este último caso, el informe se incluirá en el orden del día del siguiente pleno si la solicitud se recibe antes de su convocatoria y, si no, en el posterior.
    \item Todas aquellas que le correspondan conforme a este Reglamento y su normativa de desarrollo.
  \end{enumerate}
\end{enumerate}

\capitulo{CAPÍTULO CUARTO}{De las delegaciones y subdelegaciones de grupos de docencia de la Facultad de Ciencias}

\articulo{12}{Funciones}
\begin{enumerate}
\item Son funciones de las delegaciones y subdelegaciones de grupo:
  \begin{enumerate}
    \item Ostentar la representación del grupo de docencia o clase.
    \item Coordinarse con el personal docente y el personal técnico, de gestión y de administración y servicios, así como con los restantes representantes de estudiantes de la Facultad en defensa de los intereses de sus representados.
    \item Recabar y transmitir información en las sesiones del Pleno y a través de los representantes en el Consejo de Facultad sobre los distintos aspectos de la vida universitaria.
    \item Informar adecuadamente a sus representados acerca de todas aquellas cuestiones que afecten al desarrollo de la docencia, dando cuenta de cualquier incidencia a cualquier miembro de la Comisión Permanente.
    \item Todas aquellas que le sean atribuidas por el presente Reglamento y la legislación vigente en la materia.
  \end{enumerate}
\end{enumerate}

\articulo{13}{Normas para su elección}
\begin{enumerate}
\item Cada grupo de docencia contará con una delegación y una subdelegación, que serán elegidas anualmente, durante el primer trimestre del curso académico, mediante sufragio libre, universal, igual, secreto y directo, en un proceso en el que se asegurará en todo momento el respeto a los principios democráticos y de equidad, justicia y proporcionalidad.
\item Las elecciones serán convocadas y organizadas por el órgano competente en materia electoral de la Facultad (en adelante, Comisión Electoral) a propuesta de la presidencia del Consejo de Estudiantes. Si la presidencia no formulara la propuesta durante el primer mes del curso académico, la Comisión Electoral podrá convocarlas de oficio.
\item La determinación y, en su caso, agrupación de los grupos de docencia, de asignaturas obligatorias y de asignaturas optativas, a estos únicos efectos, deberá ser acordada por el Consejo de Facultad a propuesta del Consejo de Estudiantes.
\item El derecho de sufragio activo y pasivo corresponde en cada grupo de docencia a los estudiantes matriculados en el mismo.
\item Para optar a la delegación o subdelegación de grupo de docencia deberá presentarse la candidatura en tiempo y forma. La candidatura puede ser de un colectivo de representación estudiantil, reconocido y registrado en la Universidad, o de carácter independiente. Si se estima preciso, la Comisión Electoral propondrá la asignación a las candidaturas de los recursos necesarios para la realización de la campaña electoral.
\item La Comisión Electoral establecerá el número de mesas electorales y nombrará a sus miembros. Cada mesa estará constituida por tres estudiantes elegidos por sorteo, que no podrán en ningún caso ser candidatos en las elecciones. Además, la Comisión Electoral hará pública la convocatoria de elecciones, que deberá contener, al menos, los siguientes extremos:
  \begin{enumerate}
    \item Definición del Censo de referencia, fecha de publicación, plazo de reclamación y fecha de proclamación definitiva del Censo. En cualquier caso, la proclamación definitiva del Censo será anterior, al menos en un plazo de 24 horas, a la proclamación definitiva de candidaturas.
    \item Requisitos formales para la presentación de candidaturas.
    \item Fecha límite de presentación de candidaturas, plazos de reclamación y fecha de proclamación definitiva de candidaturas.
  \end{enumerate}
\end{enumerate}
Se dará la publicidad adecuada a la convocatoria que, en todo caso, deberá exponerse en los tablones electrónicos oficiales y enviarse a través de medios oficiales de comunicación electrónica, como la lista de correos de estudiantes de la Facultad.

\begin{enumerate}[resume]
\item Resultará elegida para la delegación la candidatura más votada y, para la subdelegación, la segunda más votada. Los empates se resolverán por sorteo, salvo acuerdo entre las personas empatadas. Si la delegación quedara vacante, la asumirá la subdelegación.
\end{enumerate}

\capitulo{CAPÍTULO QUINTO}{De los representantes designados por el Consejo de Estudiantes de la Facultad de Ciencias}

\articulo{14}{De los representantes de estudiantes en comisiones de la Facultad de Ciencias}
\begin{enumerate}
\item Los representantes de estudiantes en las comisiones de la Facultad de Ciencias serán elegidos por el Pleno de entre sus miembros, en el número que fije la normativa de cada comisión.
\item Cuando haya que elegir más de un puesto, cada miembro podrá votar a tantas candidaturas como puestos a cubrir menos uno, con un mínimo de una. Resultarán elegidas las candidaturas más votadas y las siguientes quedarán como suplentes, por orden de votos.
\item Cuando la comisión, el cargo o la representación afecte a una sola titulación, solo podrán ser candidatos y votar los miembros del Pleno matriculados en ella.
\end{enumerate}

\articulo{15}{De los representantes de estudiantes en asociaciones sectoriales}
\begin{enumerate}
\item El Pleno podrá elegir, conforme al artículo 14, representantes de estudiantes de la Facultad de Ciencias en asociaciones estudiantiles sectoriales.
\item Se fomentará la presencia de representantes en todas aquellas asociaciones sectoriales que reúnan a estudiantes de alguna de las titulaciones o ámbitos de conocimiento que se impartan en la Facultad de Ciencias.
\item La Comisión Permanente podrá financiar su participación en función de las necesidades y de la disponibilidad presupuestaria.
\end{enumerate}

\articulo{16}{De las delegaciones de titulación}
\begin{enumerate}
\item La delegación de titulación es una representación designada de estudiantes.
\item Habrá una delegación de titulación por cada titulación de grado que se imparta en la Facultad de Ciencias.
\item Podrá ser candidato cualquier miembro del Pleno matriculado en la titulación en alguna asignatura de tercer curso o superior. La elección se realizará conforme al artículo 14. Si no hubiera candidaturas, la delegación de titulación corresponderá provisionalmente, hasta que el Pleno la elija, a quien ostente la delegación de grupo del último curso; si hubiera varios grupos, la ejercerán conjuntamente, con un único voto en cualquier órgano, que acordarán entre ellas.
\item Las funciones de la delegación de titulación serán:
  \begin{enumerate}
    \item Coordinar a las delegaciones de grupo de la titulación cuando haya problemas que afecten a varios cursos.
    \item Coordinar las quejas específicas de toda la titulación.
    \item Actuar de enlace entre la coordinación de la titulación y las delegaciones de grupo de la titulación.
    \item Apoyar con su experiencia y conocimientos a las delegaciones de grupo de la titulación.
    \item Colaborar de manera directa con la coordinación de la titulación en la elaboración del calendario de exámenes del curso siguiente.
    \item Todas aquellas que le correspondan conforme a este Reglamento y su normativa de desarrollo.
  \end{enumerate}
\end{enumerate}

\capitulo{CAPÍTULO SEXTO}{Del logotipo del Consejo de Estudiantes}

\articulo{17}{Logotipo}
\begin{enumerate}
\item El logotipo del Consejo de Estudiantes de la Facultad de Ciencias es el que figura en el Anexo de este Reglamento, en su versión completa y en su versión reducida.
\item El logotipo se utilizará conforme a las normas técnicas del Manual de identidad corporativa de la Universidad de Zaragoza. La versión reducida se reservará para los usos en los que la versión completa no resulte legible, como iconos y perfiles en redes sociales.
\end{enumerate}

\disposicion{Disposición transitoria única}
Hasta la sesión constitutiva del Pleno siguiente a la entrada en vigor de este Reglamento, quienes integren el órgano ejecutivo de la Delegación de Estudiantes continuarán en funciones como Comisión Permanente, en los cargos equivalentes previstos en este Reglamento, y los representantes designados conservarán su condición.


\disposicion{Disposición derogatoria única}
Queda derogado el Reglamento de la Delegación de Alumnos de la Facultad de Ciencias, en su modificación de junio de 2015, así como cuantas normas de igual o inferior rango se opongan a lo dispuesto en este Reglamento.


\disposicion{Disposición final primera. Reforma del presente Reglamento}
La iniciativa de reforma corresponde a la Comisión Permanente o a una quinta parte de los miembros del Pleno, y el texto propuesto deberá acompañar a la convocatoria. La propuesta de reforma parcial o total de este Reglamento deberá ser aprobada por mayoría de dos tercios de los miembros presentes del Pleno, en una sesión a la que asista al menos un tercio de sus miembros, y ratificada posteriormente por el Consejo de Facultad.


\disposicion{Disposición final segunda. Derecho supletorio}
En lo no previsto en este Reglamento se aplicarán los Estatutos de la Universidad de Zaragoza, el reglamento de la Facultad de Ciencias y la Ley 40/2015, de 1 de octubre, de Régimen Jurídico del Sector Público.


\disposicion{Disposición final tercera. Entrada en vigor}
El presente Reglamento entrará en vigor al día siguiente de su ratificación por el Consejo de Facultad.


\clearpage
\disposicion{Anexo. Logotipo del Consejo de Estudiantes}

Versión completa:

\begin{center}
\includegraphics[width=110mm]{logo_consejo.pdf}
\end{center}

\vspace{1cm}
Versión reducida:

\begin{center}
\includegraphics[width=55mm]{logo_consejo_mini.pdf}
\end{center}

\vspace{1cm}
Colores: azul \#122562 y rojo \#EE2722. Tipografía: Myriad Pro.

\end{document}
